\chapter{LSTM-Autoencoder Model Architecture}
\label{ch:model_architecture}

\section{Recurrent Neural Networks and Their Limitations}
\label{sec:rnn}

Recurrent Neural Networks (RNNs) are a foundational class of neural architectures
designed to process sequential data \cite{rumelhart1986learning, goodfellow2016deep}.
Unlike feedforward networks, RNNs introduce recurrent connections that allow
information to persist across time steps. At each step $t$, the network combines the
current input $x_t$ with the previous hidden state $h_{t-1}$ to compute:

\begin{equation}
    h_t = \sigma(W_{xh} x_t + W_{hh} h_{t-1} + b_h)
    \label{eq:rnn_hidden}
\end{equation}

\noindent where $W_{xh}$ and $W_{hh}$ are learnable weight matrices, $b_h$ is a bias
vector, and $\sigma$ is a nonlinear activation (typically $\tanh$).

Despite their expressiveness, vanilla RNNs suffer from well-known training instabilities
when learning long-range dependencies \cite{bengio1994learning, pascanu2013difficulty}.
During Backpropagation Through Time (BPTT), repeated multiplication of weight matrices
across many steps causes gradients to either vanish exponentially or grow without bound.
The vanishing gradient problem is particularly harmful, as the error signal decays
before it can influence earlier states. Gradient clipping \cite{pascanu2013difficulty}
mitigates exploding gradients, but the vanishing gradient problem fundamentally
motivates the use of gated architectures such as Long Short-Term Memory networks.

\section{Long Short-Term Memory (LSTM)}
\label{sec:lstm}

\subsection{LSTM Architecture}
\label{subsec:lstm_architecture}

The LSTM network, introduced by Hochreiter and Schmidhuber
\cite{hochreiter1997long, graves2012supervised}, addresses the vanishing gradient
problem through a gating mechanism that regulates information flow within the recurrent
cell. An LSTM cell maintains a cell state $C_t$ (long-term memory) and a hidden state
$h_t$ (working memory), governed by three multiplicative gates:

\begin{itemize}
    \item \textbf{Forget Gate} ($f_t$): Determines what to discard from the previous
    cell state:
    \begin{equation}
        f_t = \sigma(W_f \cdot [h_{t-1}, x_t] + b_f)
        \label{eq:forget_gate}
    \end{equation}

    \item \textbf{Input Gate} ($i_t$): Controls what new information is written, using
    a candidate cell state $\tilde{C}_t$:
    \begin{equation}
        i_t = \sigma(W_i \cdot [h_{t-1}, x_t] + b_i)
        \label{eq:input_gate}
    \end{equation}
    \begin{equation}
        \tilde{C}_t = \tanh(W_C \cdot [h_{t-1}, x_t] + b_C)
        \label{eq:candidate_cell}
    \end{equation}

    \item \textbf{Output Gate} ($o_t$): Regulates which parts of the cell state are
    exposed as the hidden state:
    \begin{equation}
        o_t = \sigma(W_o \cdot [h_{t-1}, x_t] + b_o)
        \label{eq:output_gate}
    \end{equation}
\end{itemize}

\noindent The cell and hidden states are then updated as:

\begin{equation}
    C_t = f_t \odot C_{t-1} + i_t \odot \tilde{C}_t
    \label{eq:cell_state}
\end{equation}

\begin{equation}
    h_t = o_t \odot \tanh(C_t)
    \label{eq:hidden_state}
\end{equation}

\noindent where $\sigma$ is the sigmoid function and $\odot$ denotes element-wise
multiplication. The additive cell-state update provides a gradient highway that
preserves information across long time horizons \cite{gers2000learning}, making LSTMs
significantly more stable than vanilla RNNs for learning long-range dependencies.

\subsection{LSTM for Anomaly Detection}
\label{subsec:lstm_anomaly}

In time-series anomaly detection, LSTM-based models are trained on data representative
of normal system behavior \cite{malhotra2015lstm, hundman2018detecting,
ergen2017unsupervised}. The model learns temporal regularities and multivariate
relationships present in nominal operational patterns. During inference, sequences that
deviate from the learned representation of normality typically produce higher
reconstruction or prediction errors, enabling unsupervised fault detection without
requiring labeled anomaly examples.

\section{Autoencoders}
\label{sec:autoencoders}

\subsection{Autoencoder Architecture}
\label{subsec:ae_architecture}

An autoencoder is an unsupervised neural network that learns to encode its input into a
lower-dimensional latent representation and subsequently reconstruct the original input
\cite{hinton2006reducing, goodfellow2016deep}. The architecture consists of two
complementary sub-networks:

\begin{itemize}
    \item \textbf{Encoder} $f_\theta$: Maps the input $x \in \mathbb{R}^d$ to a
    latent vector $z \in \mathbb{R}^k$, where $k < d$.

    \item \textbf{Decoder} $g_\phi$: Reconstructs the input from the latent
    representation, producing $\hat{x} = g_\phi(z) \approx x$.
\end{itemize}

The model is trained end-to-end to minimize reconstruction error, typically using Mean
Squared Error (MSE):

\begin{equation}
    \mathcal{L}_{\text{MSE}} = \frac{1}{N} \sum_{i=1}^{N} \|x_i - \hat{x}_i\|^2
    \label{eq:mse_loss}
\end{equation}

The bottleneck imposed by the latent dimension $k$ forces the model to learn a
compressed representation of the input distribution rather than merely replicating the
raw input.

\subsection{Why Autoencoders for Anomaly Detection}
\label{subsec:ae_why}

Autoencoders are widely used for anomaly detection because they can learn a
representation of normal data without requiring explicit anomaly labels
\cite{chalapathy2019deep, pang2021deep}. Their suitability stems from several
properties:

\begin{itemize}
    \item \textbf{Unsupervised Learning}: They can be trained using only data
    representative of normal behavior, which is advantageous where anomalous examples
    are rare or difficult to label.

    \item \textbf{Multivariate Dependency Modeling}: They learn cross-feature
    relationships and compressed representations of high-dimensional data.

    \item \textbf{Reconstruction Error as Anomaly Score}: Since the model is optimized
    to reconstruct normal patterns, anomalous inputs that fall outside the learned
    distribution yield elevated reconstruction errors.
\end{itemize}

Compared to Principal Component Analysis (PCA)---a linear analog of the
autoencoder---neural autoencoders capture nonlinear feature interactions, making them
more expressive for complex telemetry signals \cite{jolliffe2016pca, chalapathy2019deep}.
The autoencoder paradigm also extends naturally to sequential variants (e.g., LSTM-AE),
whereas PCA offers no inherent temporal modeling.

\section{LSTM-Autoencoder Architecture}
\label{sec:lstm_ae}

\subsection{Combined Architecture}
\label{subsec:lstm_ae_architecture}

The LSTM-Autoencoder (LSTM-AE) integrates the temporal modeling capacity of LSTMs with
the unsupervised representation learning of autoencoders
\cite{malhotra2015lstm, zhou2017anomaly}, yielding a model well-suited to multivariate
time-series anomaly detection. The processing pipeline is:

\begin{equation}
    X \xrightarrow{\text{Encoder LSTM}} z \xrightarrow{\text{Repeat}}
    z^{(1..T)} \xrightarrow{\text{Decoder LSTM}} H
    \xrightarrow{\text{Linear}} \hat{X}
    \label{eq:lstm_ae_pipeline}
\end{equation}

\noindent where $X \in \mathbb{R}^{T \times d}$ is the input sequence of $T$ time
steps with $d$ features, $z \in \mathbb{R}^{k}$ is the compressed latent vector, and
$\hat{X} \in \mathbb{R}^{T \times d}$ is the reconstructed sequence.

\subsection{Encoder LSTM}
\label{subsec:encoder_lstm}

The encoder LSTM processes the input sequence $X = (x_1, x_2, \ldots, x_T)$
sequentially and produces a latent representation from its final hidden state:

\begin{equation}
    z = h_T \in \mathbb{R}^{k}
\end{equation}

This representation is intended to capture the most salient temporal and multivariate
characteristics of the input window.

\subsection{Latent Space}
\label{subsec:latent_space}

The latent dimensionality $k$ is a key hyperparameter governing the degree of
compression. A latent space that is too large may reconstruct even anomalous inputs
faithfully, reducing anomaly separability. One that is too small may fail to preserve
normal structure, leading to uniformly poor reconstruction. Its optimal value depends
on data complexity and the required anomaly sensitivity.

\subsection{Decoder LSTM}
\label{subsec:decoder_lstm}

The decoder LSTM reconstructs the temporal sequence from $z$ by repeating the latent
vector across the time dimension:

\begin{equation}
    z^{(1..T)} = (z, z, \ldots, z) \in \mathbb{R}^{T \times k}
\end{equation}

\noindent This repeated sequence is processed by the decoder LSTM, producing a sequence
of hidden states $H = (h_1^{\text{dec}}, \ldots, h_T^{\text{dec}})$ corresponding to
the reconstructed temporal structure.

\subsection{Linear Output Layer}
\label{subsec:linear_layer}

A shared linear transformation projects each decoder hidden state back into the
original feature space:

\begin{equation}
    \hat{x}_t = W_{\text{out}} h_t^{\text{dec}} + b_{\text{out}}
    \label{eq:linear_output}
\end{equation}

\noindent yielding the final reconstruction $\hat{X} = (\hat{x}_1, \ldots, \hat{x}_T)$.
The model is trained end-to-end by minimizing the MSE between $X$ and $\hat{X}$.

The computational footprint of the LSTM-AE is governed by input dimensionality, hidden
size, number of recurrent layers, and sequence length. In latency-sensitive or
resource-constrained deployments, compact architectures that balance representational
power with feasible inference cost are preferred.

\section{Bidirectional LSTM}
\label{sec:bilstm}

\subsection{Architecture}
\label{subsec:bilstm_arch}

A Bidirectional LSTM (BiLSTM) extends the standard LSTM by processing the input
sequence in both forward ($t = 1 \rightarrow T$) and backward ($t = T \rightarrow 1$)
directions using two separate recurrent layers \cite{graves2005framewise,
schuster1997bidirectional}. The hidden states from both passes are concatenated at each
time step:

\begin{equation}
    h_t^{\text{bi}} = [\overrightarrow{h_t} \| \overleftarrow{h_t}]
    \in \mathbb{R}^{2k}
    \label{eq:bilstm}
\end{equation}

This dual-direction processing yields richer temporal representations by incorporating
both past and future context for each position in the sequence.

\subsection{Advantages and Limitations for Edge Deployment}
\label{subsec:bilstm_tradeoffs}

Bidirectional LSTMs offer greater representational capacity, but introduce practical
trade-offs relevant to operational deployments:

\begin{itemize}
    \item \textbf{Increased Computational Cost}: Bidirectional processing requires more
    parameters and more computation than a unidirectional architecture.

    \item \textbf{Higher Memory Usage}: Maintaining forward and backward hidden states
    increases runtime memory requirements.

    \item \textbf{Full Sequence Requirement}: Backward processing requires the complete
    input sequence to be available before inference can be performed, which may limit
    suitability for low-latency streaming scenarios.
\end{itemize}

The choice between unidirectional and bidirectional architectures therefore depends on
the trade-off between representational richness and deployment constraints such as
latency, memory, and online processing requirements.

\subsection{Model Configuration}
\label{subsec:model_config}

Table~\ref{tab:model_config} summarizes the architecture parameters used across
the experimental evaluation. The selected configuration (Config~C) balances
representational capacity with edge deployment constraints.

\begin{table}[H]
    \centering
    \small
    \begin{tabular}{|l|c|l|}
        \hline
        \textbf{Parameter} & \textbf{Value} & \textbf{Rationale} \\
        \hline
        Input dimension ($d$) & 10 & Ten telemetry features \\
        \hline
        Sequence length ($T$) & 50 & $\approx$8 min at 10\,s intervals \\
        \hline
        Encoder hidden ($k$) & 128 & Best F1 in sweep (Config~C) \\
        \hline
        Decoder hidden & 128 & Symmetric to encoder \\
        \hline
        LSTM layers & 2 & Depth--accuracy trade-off \\
        \hline
        Dropout & 0.2 & Regularization for edge data \\
        \hline
        Total parameters & $\approx$410\,K & Fits in $<$2\,MB \\
        \hline
    \end{tabular}
    \caption{LSTM-AE architecture configuration}
    \label{tab:model_config}
\end{table}

\section{Training Procedure}
\label{sec:training_procedure}

The LSTM-AE is trained in an unsupervised fashion on sequences of normal
system behavior. The objective is to minimize the MSE loss
(Eq.~\ref{eq:mse_loss}) between input windows and their reconstructions.

\textbf{Optimization}: Adam \cite{kingma2015adam} with learning rate
$10^{-3}$, $\beta_1 = 0.9$, $\beta_2 = 0.999$. Batch size~64, trained
for 50~epochs per local round (5~epochs in FL mode).

\textbf{Threshold Calibration}: After training, reconstruction errors
are computed over the training set. The anomaly threshold is initialized
at the 95th percentile of training errors, serving as the starting point
for Q-learning adaptation (Chapter~\ref{ch:adaptive_thresholding}).

\textbf{Quantization}: Dynamic INT8 quantization is applied to LSTM and
Linear layers post-training via \texttt{torch.quantization.quantize\_dynamic},
reducing model size and inference latency without measurable accuracy
degradation.
